% Options for packages loaded elsewhere
% Options for packages loaded elsewhere
\PassOptionsToPackage{unicode}{hyperref}
\PassOptionsToPackage{hyphens}{url}
\PassOptionsToPackage{dvipsnames,svgnames,x11names}{xcolor}
%
\documentclass[
  english,
  letterpaper,
  DIV=11,
  numbers=noendperiod]{scrreprt}
\usepackage{xcolor}
\usepackage{amsmath,amssymb}
\setcounter{secnumdepth}{5}
\usepackage{iftex}
\ifPDFTeX
  \usepackage[T1]{fontenc}
  \usepackage[utf8]{inputenc}
  \usepackage{textcomp} % provide euro and other symbols
\else % if luatex or xetex
  \usepackage{unicode-math} % this also loads fontspec
  \defaultfontfeatures{Scale=MatchLowercase}
  \defaultfontfeatures[\rmfamily]{Ligatures=TeX,Scale=1}
\fi
\usepackage{lmodern}
\ifPDFTeX\else
  % xetex/luatex font selection
\fi
% Use upquote if available, for straight quotes in verbatim environments
\IfFileExists{upquote.sty}{\usepackage{upquote}}{}
\IfFileExists{microtype.sty}{% use microtype if available
  \usepackage[]{microtype}
  \UseMicrotypeSet[protrusion]{basicmath} % disable protrusion for tt fonts
}{}
\makeatletter
\@ifundefined{KOMAClassName}{% if non-KOMA class
  \IfFileExists{parskip.sty}{%
    \usepackage{parskip}
  }{% else
    \setlength{\parindent}{0pt}
    \setlength{\parskip}{6pt plus 2pt minus 1pt}}
}{% if KOMA class
  \KOMAoptions{parskip=half}}
\makeatother
% Make \paragraph and \subparagraph free-standing
\makeatletter
\ifx\paragraph\undefined\else
  \let\oldparagraph\paragraph
  \renewcommand{\paragraph}{
    \@ifstar
      \xxxParagraphStar
      \xxxParagraphNoStar
  }
  \newcommand{\xxxParagraphStar}[1]{\oldparagraph*{#1}\mbox{}}
  \newcommand{\xxxParagraphNoStar}[1]{\oldparagraph{#1}\mbox{}}
\fi
\ifx\subparagraph\undefined\else
  \let\oldsubparagraph\subparagraph
  \renewcommand{\subparagraph}{
    \@ifstar
      \xxxSubParagraphStar
      \xxxSubParagraphNoStar
  }
  \newcommand{\xxxSubParagraphStar}[1]{\oldsubparagraph*{#1}\mbox{}}
  \newcommand{\xxxSubParagraphNoStar}[1]{\oldsubparagraph{#1}\mbox{}}
\fi
\makeatother


\usepackage{longtable,booktabs,array}
\newcounter{none} % for unnumbered tables
\usepackage{calc} % for calculating minipage widths
% Correct order of tables after \paragraph or \subparagraph
\usepackage{etoolbox}
\makeatletter
\patchcmd\longtable{\par}{\if@noskipsec\mbox{}\fi\par}{}{}
\makeatother
% Allow footnotes in longtable head/foot
\IfFileExists{footnotehyper.sty}{\usepackage{footnotehyper}}{\usepackage{footnote}}
\makesavenoteenv{longtable}
\usepackage{graphicx}
\makeatletter
\newsavebox\pandoc@box
\newcommand*\pandocbounded[1]{% scales image to fit in text height/width
  \sbox\pandoc@box{#1}%
  \Gscale@div\@tempa{\textheight}{\dimexpr\ht\pandoc@box+\dp\pandoc@box\relax}%
  \Gscale@div\@tempb{\linewidth}{\wd\pandoc@box}%
  \ifdim\@tempb\p@<\@tempa\p@\let\@tempa\@tempb\fi% select the smaller of both
  \ifdim\@tempa\p@<\p@\scalebox{\@tempa}{\usebox\pandoc@box}%
  \else\usebox{\pandoc@box}%
  \fi%
}
% Set default figure placement to htbp
\def\fps@figure{htbp}
\makeatother



\ifLuaTeX
\usepackage[bidi=basic,shorthands=off]{babel}
\else
\usepackage[bidi=default,shorthands=off]{babel}
\fi
\ifLuaTeX
  \usepackage{selnolig} % disable illegal ligatures
\fi


\setlength{\emergencystretch}{3em} % prevent overfull lines

\providecommand{\tightlist}{%
  \setlength{\itemsep}{0pt}\setlength{\parskip}{0pt}}



 


\KOMAoption{captions}{tableheading}
\makeatletter
\@ifpackageloaded{bookmark}{}{\usepackage{bookmark}}
\makeatother
\makeatletter
\@ifpackageloaded{caption}{}{\usepackage{caption}}
\AtBeginDocument{%
\ifdefined\contentsname
  \renewcommand*\contentsname{Table of contents}
\else
  \newcommand\contentsname{Table of contents}
\fi
\ifdefined\listfigurename
  \renewcommand*\listfigurename{List of Figures}
\else
  \newcommand\listfigurename{List of Figures}
\fi
\ifdefined\listtablename
  \renewcommand*\listtablename{List of Tables}
\else
  \newcommand\listtablename{List of Tables}
\fi
\ifdefined\figurename
  \renewcommand*\figurename{Figure}
\else
  \newcommand\figurename{Figure}
\fi
\ifdefined\tablename
  \renewcommand*\tablename{Table}
\else
  \newcommand\tablename{Table}
\fi
}
\@ifpackageloaded{float}{}{\usepackage{float}}
\floatstyle{ruled}
\@ifundefined{c@chapter}{\newfloat{codelisting}{h}{lop}}{\newfloat{codelisting}{h}{lop}[chapter]}
\floatname{codelisting}{Listing}
\newcommand*\listoflistings{\listof{codelisting}{List of Listings}}
\makeatother
\makeatletter
\makeatother
\makeatletter
\@ifpackageloaded{caption}{}{\usepackage{caption}}
\@ifpackageloaded{subcaption}{}{\usepackage{subcaption}}
\makeatother
\usepackage{bookmark}
\IfFileExists{xurl.sty}{\usepackage{xurl}}{} % add URL line breaks if available
\urlstyle{same}
\hypersetup{
  pdftitle={PGM: Practical Growth Modeling for jamovi},
  pdfauthor={André L. R. Pessina and Manoel Garcia Neto},
  pdflang={en},
  colorlinks=true,
  linkcolor={blue},
  filecolor={Maroon},
  citecolor={Blue},
  urlcolor={Blue},
  pdfcreator={LaTeX via pandoc}}


\title{PGM: Practical Growth Modeling for jamovi}
\usepackage{etoolbox}
\makeatletter
\providecommand{\subtitle}[1]{% add subtitle to \maketitle
  \apptocmd{\@title}{\par {\large #1 \par}}{}{}
}
\makeatother
\subtitle{Methods and User Guide}
\author{André L. R. Pessina and Manoel Garcia Neto}
\date{}
\begin{document}
\maketitle

\renewcommand*\contentsname{Table of contents}
{
\setcounter{tocdepth}{2}
\tableofcontents
}

\bookmarksetup{startatroot}

\chapter{Index}\label{index}

\bookmarksetup{startatroot}

\chapter{PGM}\label{pgm}

\textbf{PGM --- Practical Growth Models} is a jamovi module for fitting,
evaluating, and interpreting nonlinear growth curves.

The module provides tools for:

\begin{itemize}
\tightlist
\item
  fitting parametric growth models;
\item
  modeling heteroscedastic residual variance;
\item
  fitting nonlinear mixed-effects models for longitudinal data;
\item
  modeling residual autocorrelation;
\item
  evaluating model fit and prediction error;
\item
  estimating biologically and geometrically relevant growth points;
\item
  visualizing fitted curves, derivatives, and residual diagnostics.
\end{itemize}

\section{Scope}\label{scope}

PGM is intended for researchers analyzing longitudinal or
cross-sectional growth data in animal science, plant science,
microbiology, ecology, developmental biology, and related fields.

The documentation is divided into four main parts:

\begin{enumerate}
\def\labelenumi{\arabic{enumi}.}
\tightlist
\item
  \textbf{Getting started}, covering installation, data organization,
  and a first analysis;
\item
  \textbf{Methods}, describing the models and estimation procedures
  implemented in PGM;
\item
  \textbf{Results}, explaining the output tables, diagnostics, and
  derived growth descriptors;
\item
  \textbf{Academic}, providing reporting guidance and limitations.
\end{enumerate}

\section{Important considerations}\label{important-considerations}

PGM assists with model estimation and interpretation but does not
replace statistical judgment. Model selection should not be based on a
single goodness-of-fit statistic. The adequacy of the mean function,
residual variance, dependence structure, parameter plausibility, and
residual diagnostics should be considered jointly.

\section{Acknowledgments and Contact}\label{acknowledgments-and-contact}

This project is being developed at the School of Veterinary Medicine,
São Paulo State University (UNESP), Araçatuba, São Paulo, Brazil.

We extend our special thanks to the jamovi development team,
particularly Jonathon Love, for his support, as well as to all users and
testers who have contributed to the development of this project.

Comments and feedback are welcome. Please contact us if you identify any
errors, have suggestions or questions, or require methodological support
for academic research. The development team can be reached at
andre.pessina@unesp.br

Citation information for the current version is provided in the chapter
on academic reporting. In addition to citing PGM, we encourage authors
to cite the R Project and the jamovi project. We also recommend citing
the packages and references used in the module's analyses.

\bookmarksetup{startatroot}

\chapter{Introduction}\label{introduction}

\bookmarksetup{startatroot}

\chapter{Introduction}\label{introduction-1}

Growth is a dynamic process in which the size, mass, concentration, or
other measurable characteristic of a biological system changes over
time.

Sigmoidal growth models summarize an entire trajectory using a small
number of parameters related to characteristics such as:

\begin{itemize}
\tightlist
\item
  asymptotic response \texttt{A};
\item
  growth rate \texttt{K};
\item
  timing of the inflection point \texttt{Ti};
\item
  symmetry or shape of the trajectory \texttt{d}.
\end{itemize}

PGM provides a graphical interface for fitting these models in jamovi,
while retaining statistical procedures commonly implemented in R.

\section{Why use parametric growth
models?}\label{why-use-parametric-growth-models}

Parametric growth models can provide interpretable summaries of
nonlinear trajectories. Depending on the selected model, their
parameters may describe:

\begin{itemize}
\tightlist
\item
  the expected mature or asymptotic value;
\item
  the temporal scale of growth;
\item
  the maximum growth rate;
\item
  the age or time of maximum growth rate;
\item
  the shape and asymmetry of the curve.
\end{itemize}

These interpretations depend on the equation and parametrization used.
Parameters with similar names are not necessarily equivalent across
different growth models.

\section{Statistical challenges}\label{statistical-challenges}

Biological growth data frequently present characteristics that are not
adequately represented by ordinary nonlinear least squares, including:

\begin{itemize}
\tightlist
\item
  increasing or decreasing residual variability;
\item
  repeated observations from the same experimental unit;
\item
  between-individual variation in growth parameters;
\item
  serial correlation between observations;
\item
  incomplete or irregular measurement schedules.
\end{itemize}

PGM therefore includes variance functions, nonlinear mixed-effects
models, and residual correlation structures.

\section{Documentation conventions}\label{documentation-conventions}

Throughout this documentation:

\begin{itemize}
\tightlist
\item
  \(W(t)\) denotes the fitted growth function;
\item
  \(W'(t)\) denotes the instantaneous growth rate;
\item
  \(W''(t)\) denotes growth acceleration;
\item
  \(t\) denotes time or age;
\item
  \(y\) denotes the observed response.
\end{itemize}

The interpretation of each model-specific parameter is presented in the
chapter on growth models.

\bookmarksetup{startatroot}

\chapter{Installation}\label{installation}

\bookmarksetup{startatroot}

\chapter{Installation}\label{installation-1}

\section{Installing jamovi}\label{installing-jamovi}

PGM is an add-on module for jamovi. Before installing the module,
download and install the current version of jamovi from its
\href{https://www.jamovi.org/download.html}{official website}.

After installation, open jamovi and verify that the main analysis
interface is working correctly.

\section{Installing PGM from the jamovi
library}\label{installing-pgm-from-the-jamovi-library}

PGM is available in the official jamovi library:

\begin{enumerate}
\def\labelenumi{\arabic{enumi}.}
\item
  Open jamovi.
\item
  Select the Modules icon, in the upper-right of Analysis menu.
\item
  Open the jamovi module library.
\item
  Search for PGM.
\item
  Select Install.
\end{enumerate}

After installation, PGM will appear in the main analysis ribbon.

\section{Installing a development
version}\label{installing-a-development-version}

Development versions may be distributed as .jmo files on
\href{https://github.com/apessina/jamovi-PGM}{PGM GitHub repository}.

\begin{enumerate}
\def\labelenumi{\arabic{enumi}.}
\item
  To install a local module:
\item
  Open jamovi.
\item
  Select the Modules icon.
\item
  Select Sideload.
\item
  Locate the .jmo file.
\item
  Confirm the installation.
\end{enumerate}

Development versions may contain recently added features that have not
yet been included in a stable release.

\section{Updating PGM}\label{updating-pgm}

When PGM is installed through the jamovi library, available updates can
be installed through the module manager.

For reproducible academic work, record the PGM and jamovi versions.

\section{Verifying the installation}\label{verifying-the-installation}

After installation, open a dataset containing numeric time and response
variables. The module includes example datasets for testing and learning
purposes. They are available from the jamovi menu in the upper-left
corner: \textbf{Open → Data Library → PGM}.

Next, select \textbf{Growth} from the analysis ribbon and open
\textbf{Curve Modeling}. Assign at least one response variable and one
time variable.

If the analysis interface opens and produces an initial result, the
module has been installed successfully.

\section{Reporting problems}\label{reporting-problems}

When reporting an error, include:

\begin{itemize}
\item
  the PGM version;
\item
  the jamovi version;
\item
  the operating system;
\item
  the selected model and options;
\item
  the complete error message;
\item
  a reproducible or anonymized dataset, when possible.
\end{itemize}

\bookmarksetup{startatroot}

\chapter{Data}\label{data}

\bookmarksetup{startatroot}

\chapter{Data preparation}\label{data-preparation}

\section{Example datasets}\label{example-datasets}

The module includes example datasets derived from published studies for
testing and learning purposes. References for the original data sources
are provided with each dataset.

They are available from the jamovi menu in the upper-left corner:
\textbf{Open → Data Library → PGM}.

\section{Basic structure}\label{basic-structure}

Data should be provided in long format, with one row per observation.

At minimum, the dataset must contain:

\begin{itemize}
\tightlist
\item
  one numeric response variable;
\item
  one numeric time or age variable.
\end{itemize}

For longitudinal analyses, the dataset must also contain:

\begin{itemize}
\tightlist
\item
  one identification variable indicating the experimental unit or
  subject.
\end{itemize}

\section{Cross-sectional structure}\label{cross-sectional-structure}

A dataset without an identification variable is interpreted as
containing independent observations.

{\def\LTcaptype{none} % do not increment counter
\begin{longtable}[]{@{}rr@{}}
\toprule\noalign{}
time & response \\
\midrule\noalign{}
\endhead
\bottomrule\noalign{}
\endlastfoot
1 & 12.4 \\
2 & 18.9 \\
3 & 27.5 \\
\end{longtable}
}

\section{Longitudinal structure}\label{longitudinal-structure}

For repeated measurements, each subject must have a stable identifier.

{\def\LTcaptype{none} % do not increment counter
\begin{longtable}[]{@{}lrr@{}}
\toprule\noalign{}
subject & time & response \\
\midrule\noalign{}
\endhead
\bottomrule\noalign{}
\endlastfoot
A01 & 1 & 12.4 \\
A01 & 2 & 18.9 \\
A01 & 3 & 27.5 \\
A02 & 1 & 11.8 \\
A02 & 2 & 19.7 \\
A02 & 3 & 29.1 \\
\end{longtable}
}

\section{Variable requirements}\label{variable-requirements}

\subsection{Response}\label{response}

The response must be numeric and should represent a cumulative or
monotonic biological growth process when a sigmoidal growth model is
used.

\subsection{Time}\label{time}

Time must be numeric. Examples include:

\begin{itemize}
\tightlist
\item
  age in days;
\item
  age in weeks;
\item
  hours after inoculation;
\item
  days after planting.
\end{itemize}

The unit of time determines the unit of rate parameters and derivative
quantities.

\subsection{Subject identifier}\label{subject-identifier}

The subject identifier may be numeric or categorical. It is used to
associate repeated observations with the same experimental unit.

Selecting an identifier changes the statistical model from an
\emph{ordinary or generalized nonlinear model} to a \emph{nonlinear
mixed-effects model}.

\section{Missing observations}\label{missing-observations}

Rows containing missing values in variables required by a given analysis
are excluded from that analysis.

Different response variables may therefore be fitted using different
numbers of observations.

\section{Measurement schedule}\label{measurement-schedule}

For models using discrete AR(1) residual correlation, equally spaced
observations are generally expected.

For irregularly spaced measurements, continuous-time CAR(1) correlation
is usually more appropriate.

\section{Data coverage}\label{data-coverage}

A nonlinear growth curve should ideally contain observations from the
early, rapidly growing, and late portions of the trajectory.

If the observed data do not adequately cover the upper part of the
curve, the asymptotic parameter may be weakly identified and highly
uncertain.

\bookmarksetup{startatroot}

\chapter{First analysis}\label{first-analysis}

\bookmarksetup{startatroot}

\chapter{Running an analysis}\label{running-an-analysis}

\section{Opening the analysis}\label{opening-the-analysis}

To begin:

\begin{enumerate}
\def\labelenumi{\arabic{enumi}.}
\item
  Open the dataset in jamovi.
\item
  Select PGM from the analysis ribbon.
\item
  Open the growth-curve analysis.
\item
  Assign the time variable.
\item
  Assign one or more response variables.
\item
  Add a modeling setup for each response-model combination.
\end{enumerate}

Each modeling setup defines the response variable, growth model,
residual variance structure, and residual correlation structure to be
used in that analysis.

\section{Selecting the response}\label{selecting-the-response}

The response should be numeric and should represent a process compatible
with the selected growth equation.

Examples include: body weight; organism size; microbial concentration;
plant height; cumulative production; other continuous responses measured
over time.

A sigmoidal growth model should not be selected solely because the
observations increase over time. The observed trajectory should also
exhibit a biologically and statistically plausible sigmoidal form.

\section{Selecting time}\label{selecting-time}

The time variable must be numeric.

\textbf{Its measurement unit determines the units of:}

\begin{itemize}
\item
  rate parameters;
\item
  inflection time;
\item
  critical-point locations;
\item
  derivative-based descriptors;
\item
  autocorrelation distances.
\end{itemize}

Changing time from days to weeks changes the numerical values of
time-dependent parameters.

\section{Subject identifier}\label{subject-identifier-1}

For repeated measurements, assign a subject or experimental-unit
identifier.

\textbf{When no identifier is supplied, observations are treated as
independent.}

When an identifier is supplied, PGM fits a \emph{nonlinear mixed-effects
model} in which repeated observations are grouped within subjects.

The identifier should represent the unit that generates the repeated
trajectory, such as: animal; plant; patient; experimental vessel;
microbial batch; plot.

Treatment group, sex, breed, or experimental condition should not be
used as the subject identifier unless each category genuinely represents
one repeated observational unit.

\section{Selecting a model}\label{selecting-a-model}

Choose the growth equation according to: the expected trajectory shape;
parameter interpretability; biological context; data coverage;
convergence; residual diagnostics; comparative model performance.

Different equations may fit similarly while producing parameters with
different meanings.

\textbf{\emph{The current version fits only the 4-parameter Richards,
that should cover the shapes of both Gompertz and Logistic models.}}

\section{Selecting a residual variance
structure}\label{selecting-a-residual-variance-structure}

Use a variance function when residual spread changes systematically with
the fitted response.

Available structures may include:

\begin{itemize}
\item
  constant variance;
\item
  power variance;
\item
  exponential variance;
\item
  constant-plus-power variance.
\end{itemize}

The selection should be supported by residual diagnostics and
comparative model evaluation.

Selecting residual correlation

Residual correlation structures are available for longitudinal analyses.

Use:

AR(1) when measurements are approximately equally spaced and correlation
is interpreted by measurement lag; CAR(1) when measurement intervals are
irregular and correlation depends on the actual temporal distance.

Random effects and residual autocorrelation represent different
dependence mechanisms and may be included simultaneously.

\section{Recommended workflow}\label{recommended-workflow}

A defensible analysis commonly follows this sequence:

\begin{enumerate}
\def\labelenumi{\arabic{enumi}.}
\item
  inspect the raw trajectories;
\item
  fit a biologically plausible mean curve;
\item
  assess parameter estimates and convergence;
\item
  inspect residuals versus fitted values and time;
\item
  evaluate heteroscedasticity;
\item
  account for repeated measurements using a subject identifier;
\item
  investigate residual autocorrelation;
\item
  compare candidate specifications;
\item
  verify parameter plausibility;
\item
  interpret derived growth descriptors only after accepting the
  underlying fitted model.
\end{enumerate}

\section{Multiple modeling setups}\label{multiple-modeling-setups}

PGM allows multiple response-model configurations in the same analysis.

This can be used to:

\begin{itemize}
\item
  compare candidate growth models;
\item
  compare variance structures;
\item
  compare response variables;
\item
  inspect alternative residual-correlation structures.
\end{itemize}

Comparisons should only be interpreted formally when the models use the
same response observations and compatible likelihoods.

\part{Methods}

\chapter{overview}\label{overview}

\chapter{Modeling framework}\label{modeling-framework}

PGM selects the estimation framework according to the structure of the
data and the analysis options.

\section{Estimation pathways}\label{estimation-pathways}

{\def\LTcaptype{none} % do not increment counter
\begin{longtable}[]{@{}ccl@{}}
\toprule\noalign{}
Subject identifier & Residual variance & Estimation framework \\
\midrule\noalign{}
\endhead
\bottomrule\noalign{}
\endlastfoot
No & Constant & Nonlinear least squares \\
No & Modeled & Generalized nonlinear least squares \\
Yes & Any available option & Nonlinear mixed-effects model \\
\end{longtable}
}

All estimation pathways begin with an initial nonlinear fit using the
Levenberg--Marquardt algorithm.

When a residual variance function is selected without a subject
identifier, the initial estimates are used to initialize a generalized
nonlinear least-squares model.

When a subject identifier is selected, the initial estimates are used to
initialize a nonlinear mixed-effects model.

\section{Mean, variance, and
dependence}\label{mean-variance-and-dependence}

The complete statistical specification may contain three components:

\begin{enumerate}
\def\labelenumi{\arabic{enumi}.}
\tightlist
\item
  \textbf{Mean structure:} the selected parametric growth equation;
\item
  \textbf{Variance structure:} the relationship between residual
  variance and fitted response;
\item
  \textbf{Dependence structure:} random effects and, optionally,
  residual autocorrelation.
\end{enumerate}

These components address different features of the data and should not
be treated as interchangeable.

\section{Maximum likelihood}\label{maximum-likelihood}

Nonlinear mixed-effects models in PGM are estimated using maximum
likelihood.

This facilitates comparisons between nested models with different fixed
effects, provided that the models use the same response, observations,
and likelihood framework.

\chapter{growth-models}\label{growth-models}

Growth models General formulation

A parametric growth model represents the expected response at time (t)
as

\[
y_i = W(t_i;\boldsymbol\theta) + \varepsilon_i,
\]

where:

(W(t\_i;\boldsymbol\theta)) is the selected growth function;
(\boldsymbol\theta) is the vector of model parameters; (\varepsilon\_i)
is the residual error.

The interpretation of (\boldsymbol\theta) depends on the exact equation
and parametrization.

Models with similar names may be expressed using different parameter
definitions in different publications and software implementations.

Parameter units

The units of each parameter should be determined from the model
equation.

Common interpretations include:

(A): response unit; (T\_i): time unit; rate parameters: inverse time or
a parametrization-specific rate unit; shape parameters: dimensionless.

Parameters should not be compared directly across equations unless their
mathematical definitions are equivalent.

Richards model

PGM implements the Richards parametrization described by Tjørve and
Tjørve.

The growth function is

\[
A
\left[
1+
(d-1)
\exp\left(
-\frac{K(t-T_i)}
{d^{d/(1-d)}}
\right)
\right]^{1/(1-d)}.
\]

The parameters are:

(A): asymptotic response; (K): growth-rate parameter under this
parametrization; (T\_i): time at the inflection point; (d): shape
parameter controlling curve asymmetry.

The meaning of (K) is specific to this parametrization and should not be
assumed to equal the rate parameter of another Richards equation.

Parameter constraints

PGM fits the Richards model using transformed parameters:

\[
A=\exp(\alpha),
\]

\[
K=\exp(\kappa),
\]

and

\[
d=1.05+\exp(\delta).
\]

Consequently:

(A\textgreater0); (K\textgreater0); (d\textgreater1.05).

The (d\textgreater1.05) restriction is part of the current computational
implementation and restricts the range of Richards shapes that can be
fitted.

Starting values

Starting values are generated automatically from the observed
trajectory.

The initial asymptote is based on a value above the observed maximum
response. The initial inflection time is informed by the largest
numerical growth rate, and the initial rate parameter is derived from
that rate relative to the initial asymptote.

Automatic starting values facilitate routine use but do not guarantee
convergence or parameter identifiability.

Documenting additional models

For every additional model included in PGM, this chapter should provide:

the exact implemented equation; the source or parametrization reference;
parameter definitions; parameter units; parameter constraints; the
location of the inflection point; limiting or nested cases, when
relevant; cautions about comparison with alternative parametrizations.

The equation shown in this manual must match the equation implemented by
the module, rather than a generically named version found elsewhere in
the literature.

\chapter{estimation}\label{estimation}

\chapter{Estimation}\label{estimation-1}

\section{Initial nonlinear fit}\label{initial-nonlinear-fit}

PGM first fits the selected growth equation using nonlinear least
squares and the Levenberg--Marquardt algorithm.

The initial fit serves two purposes:

\begin{enumerate}
\def\labelenumi{\arabic{enumi}.}
\tightlist
\item
  it provides estimates for an ordinary nonlinear analysis;
\item
  it supplies starting values for generalized nonlinear least-squares or
  nonlinear mixed-effects estimation.
\end{enumerate}

Starting values are generated automatically from the observed data and
the selected growth model.

\section{Parameter transformations}\label{parameter-transformations}

Some models are fitted using transformed parameters to impose parameter
constraints and improve numerical stability.

For example, a positive parameter (A) may be represented internally as

\[
A = \exp(\alpha).
\]

The reported estimate is returned to the natural parameter scale.

\section{Standard errors}\label{standard-errors}

The covariance matrix is initially obtained on the transformed
estimation scale.

For parameters reported after nonlinear back-transformation, uncertainty
is propagated to the natural scale using a first-order delta method.

If (\boldsymbol\theta) denotes the transformed parameters and
(g(\boldsymbol\theta)) their natural-scale representation, then

\[
\operatorname{Var}\!\left[g(\hat{\boldsymbol\theta})\right]
\approx
J
\operatorname{Var}(\hat{\boldsymbol\theta})
J^\mathsf{T},
\]

where (J) is the Jacobian matrix of the back-transformation.

\section{Confidence intervals}\label{confidence-intervals}

Approximate 95\% confidence intervals are first constructed on the
transformed scale and subsequently back-transformed.

These intervals depend on asymptotic normal approximations and may be
unreliable when:

\begin{itemize}
\tightlist
\item
  the sample size is small;
\item
  parameters are weakly identified;
\item
  estimates are near parameter boundaries;
\item
  the fitted model is numerically unstable.
\end{itemize}

\chapter{mixed-effects}\label{mixed-effects}

\chapter{Nonlinear mixed-effects
models}\label{nonlinear-mixed-effects-models}

\section{When mixed effects are used}\label{when-mixed-effects-are-used}

A nonlinear mixed-effects model is fitted when a subject identification
variable is supplied.

The model separates:

\begin{itemize}
\tightlist
\item
  population-level growth parameters, represented by fixed effects;
\item
  individual deviations from the population trajectory, represented by
  random effects;
\item
  within-individual residual variation.
\end{itemize}

\section{Current random-effects
structure}\label{current-random-effects-structure}

In the current implementation, the asymptotic parameter (A) contains an
individual random effect.

On the transformed scale, the subject-specific asymptote may be written
as

\[
\log(A_i) = \log(A) + b_{Ai},
\]

where

\[
b_{Ai} \sim N(0,\sigma_A^2).
\]

All remaining growth parameters are treated as fixed population
parameters.

\section{Interpretation}\label{interpretation}

The fixed asymptotic parameter represents the population-level central
trajectory on the model scale.

The random effect represents persistent between-individual variation in
the asymptotic level.

Residual error represents observation-level variation remaining after
accounting for the fixed effects, random effects, variance function, and
residual correlation structure.

\section{Conditional fitted values}\label{conditional-fitted-values}

Residual diagnostics for nonlinear mixed-effects models use conditional
fitted values. These include the estimated individual random effects.

Consequently, the diagnostic residuals describe fit to individual
trajectories rather than only to the population-average trajectory.

\chapter{variance-functions}\label{variance-functions}

\chapter{Residual variance functions}\label{residual-variance-functions}

Ordinary nonlinear least squares assumes constant residual variance:

\[
\operatorname{Var}(\varepsilon_i)=\sigma^2.
\]

This assumption is often unrealistic for biological growth data because
larger responses may exhibit greater absolute variability.

PGM provides alternative residual variance functions based on the fitted
response (\hat y\_i).

\section{Power variance function}\label{power-variance-function}

\[
\operatorname{Var}(\varepsilon_i)
=
\sigma^2 |\hat y_i|^{2\theta}.
\]

Interpretation:

\begin{itemize}
\tightlist
\item
  (\theta=0) corresponds to constant variance;
\item
  (\theta\textgreater0) indicates increasing variance with fitted
  response;
\item
  (\theta\textless0) indicates decreasing variance with fitted response.
\end{itemize}

\section{Exponential variance
function}\label{exponential-variance-function}

\[
\operatorname{Var}(\varepsilon_i)
=
\sigma^2 \exp(2\theta\hat y_i).
\]

This model represents an exponential relationship between residual
standard deviation and fitted response.

\section{Constant-plus-power variance
function}\label{constant-plus-power-variance-function}

\[
\operatorname{Var}(\varepsilon_i)
=
\sigma^2
\left(
\theta_1 + |\hat y_i|^{\theta_2}
\right)^2.
\]

This structure permits a nonzero baseline component together with a
power-related increase or decrease in residual variation.

\section{Choosing a variance
structure}\label{choosing-a-variance-structure}

A variance function should be considered when residual spread changes
systematically across fitted values or time.

Selection should consider:

\begin{itemize}
\tightlist
\item
  residual-versus-fitted plots;
\item
  normalized residuals;
\item
  information criteria;
\item
  parameter stability;
\item
  biological plausibility.
\end{itemize}

A lower AIC alone is insufficient if the fitted variance structure
produces unstable parameters or implausible trajectories.

\chapter{residual-correlation}\label{residual-correlation}

\chapter{Residual correlation}\label{residual-correlation-1}

Repeated observations from the same subject may remain correlated even
after individual random effects are included.

PGM provides first-order discrete and continuous-time correlation
structures.

\section{AR(1)}\label{ar1}

For observations separated by (h) measurement steps,

\[
\operatorname{Corr}
(\varepsilon_{i,t},\varepsilon_{i,t+h})
=
\phi^{|h|},
\]

where (-1 \textless{} \phi \textless{} 1).

AR(1) is most naturally interpreted when observations are ordered and
approximately equally spaced.

\section{CAR(1)}\label{car1}

For observations made at times (t\_\{ij\}) and (t\_\{ik\}),

\[
\operatorname{Corr}
(\varepsilon_{ij},\varepsilon_{ik})
=
\phi^{|t_{ij}-t_{ik}|}.
\]

CAR(1) uses the actual temporal distance between observations and is
therefore suitable for irregularly spaced measurements.

\section{Random effects are not
autocorrelation}\label{random-effects-are-not-autocorrelation}

Random effects and residual correlation represent different sources of
dependence.

A random effect represents persistent individual differences in the
growth trajectory.

Residual autocorrelation represents association between
observation-level departures occurring close together in time.

A model may require one, both, or neither structure.

\part{Results}

\chapter{model-evaluation}\label{model-evaluation}

Results and interpretation Estimation information

The estimation-information table describes the computational and
statistical framework used for each modeling setup.

Depending on the selected options, it may report:

nonlinear least squares; generalized nonlinear least squares; nonlinear
mixed-effects estimation; Levenberg--Marquardt initialization;
Gauss--Newton estimation; residual variance structure; residual
correlation structure; random-effects structure.

This table should be consulted before reporting the model because the
same mean equation may be fitted under different residual and dependence
assumptions.

Parameter estimates

The parameter table reports estimates on the natural parameter scale.

Available columns may include:

parameter estimate; standard error; lower confidence limit; upper
confidence limit.

Some parameters are estimated internally on a transformed scale and
subsequently returned to their natural scale.

Approximate uncertainty is propagated through the back-transformation
using a first-order delta method.

Interpreting confidence intervals

Confidence intervals quantify uncertainty under the fitted model and its
asymptotic approximation.

Wide intervals may indicate:

insufficient sample size; inadequate coverage of the growth trajectory;
strong correlation among model parameters; weak identification;
excessive model flexibility; numerical instability.

A narrow confidence interval does not establish biological validity if
the selected model is structurally inappropriate.

Goodness-of-fit measures Akaike information criterion

The Akaike information criterion is

-2\ell + 2k, \$\$

where (\ell) is the maximized log-likelihood and (k) is the number of
estimated parameters counted by the fitted likelihood.

Lower AIC values indicate better expected predictive information
performance among comparable models.

The absolute AIC value is not interpreted independently.

Corrected Akaike information criterion

For finite samples, PGM calculates

\mathrm{AIC}

\begin{itemize}
\item
  \frac{2k(k+1)}

  \{n-k-1\}. \$\$
\end{itemize}

AICc is unavailable when the sample size is not sufficiently larger than
the parameter count.

Bayesian information criterion

The Bayesian information criterion is

-2\ell+k\log(n). \$\$

BIC generally penalizes model complexity more strongly than AIC as the
sample size increases.

Descriptive coefficient of determination

PGM calculates

1- \frac{\sum_i(y_i-\hat y_i)^2} \{\sum\_i(y\_i-\bar y)\^{}2\}. \$\$

This is a descriptive residual-based statistic on the response scale.

For nonlinear mixed-effects models, it should not be described as the
marginal or conditional (R\^{}2) commonly defined specifically for mixed
models.

Adjusted descriptive coefficient of determination

PGM calculates an adjusted form based on residual degrees of freedom:

1- \frac{\mathrm{SSE}/(n-p)} \{\mathrm{SST}/(n-1)\}. \$\$

Its interpretation remains descriptive.

Error metrics Root mean squared error

\sqrt{
\frac{1}{n}
\sum_i
(y_i-\hat y_i)^2
}. \$\$

RMSE uses the response unit and is relatively sensitive to large
residuals.

Relative RMSE

\frac{\mathrm{RMSE}}{\bar y}. \$\$

RRMSE is dimensionless but may behave poorly when the response mean is
close to zero.

Mean absolute error

\frac{1}{n}

\sum\_i \textbar y\_i-\hat y\_i\textbar. \$\$

Median absolute error

\operatorname{median}

\left( \textbar y\_i-\hat y\_i\textbar{} \right). \$\$

MedAE is less sensitive to isolated large residuals than RMSE or MAE.

Symmetric mean absolute percentage error

PGM calculates

\frac{100}{n}

\sum\_i \frac{2|y_i-\hat y_i|}
\{\textbar y\_i\textbar+\textbar{}\hat y\_i\textbar\}. \$\$

A small numerical safeguard is used when both the observed and fitted
values are close to zero.

Comparing models

AIC, AICc, and BIC comparisons require:

the same response variable; the same response scale; the same
observations; compatible likelihood definitions; the same treatment of
missing values.

RMSE, MAE, and related error metrics should also be compared only across
the same observations.

Model selection should combine statistical performance, residual
diagnostics, parameter plausibility, convergence, and scientific
interpretability.

\chapter{diagnostics}\label{diagnostics}

\chapter{Residual diagnostics}\label{residual-diagnostics}

PGM provides a four-panel residual diagnostic display.

\section{Residuals versus fitted
values}\label{residuals-versus-fitted-values}

This plot is used to investigate:

\begin{itemize}
\tightlist
\item
  nonconstant residual variance;
\item
  systematic curvature;
\item
  regions of poor fit;
\item
  influential or extreme observations.
\end{itemize}

A satisfactory pattern should be centered near zero without pronounced
systematic structure.

\section{Residuals versus time}\label{residuals-versus-time}

This plot is used to investigate:

\begin{itemize}
\tightlist
\item
  temporal trends not represented by the growth function;
\item
  early- or late-stage model inadequacy;
\item
  changes in residual variability over time;
\item
  local clusters of positive or negative residuals.
\end{itemize}

\section{Normal Q--Q plot}\label{normal-qq-plot}

The Q--Q plot compares the empirical residual distribution with the
normal distribution.

Large systematic departures from the diagonal may indicate skewness,
heavy tails, outliers, or model misspecification.

Minor departures should not be interpreted mechanically, particularly
with large datasets.

\section{Residual autocorrelation}\label{residual-autocorrelation}

The autocorrelation plot assesses whether residuals remain associated
across neighboring measurement occasions.

For longitudinal models, residual autocorrelation is evaluated within
subjects.

Persistent residual autocorrelation may indicate that the selected
random effects or residual correlation structure is insufficient.

\section{Residual scale}\label{residual-scale}

For generalized nonlinear and nonlinear mixed-effects models, PGM uses
normalized residuals that account for the fitted variance and
correlation structure.

For ordinary nonlinear least squares, raw residuals are standardized by
the estimated residual standard deviation.

\chapter{critical-points}\label{critical-points}

Growth descriptors Derivative notation

Let (W(t)) denote the fitted growth function.

PGM may calculate:

\[
W'(t),
\]

the instantaneous growth rate, and

\[
W''(t),
\]

the growth acceleration.

Higher-order derivatives are used for selected geometric descriptors.

These quantities are properties of the fitted curve and inherit any
misspecification or uncertainty in the underlying growth model.

Inflection point

The inflection point is identified from a sign change in the second
derivative:

\[
W''(t)=0.
\]

At the inflection point of a standard sigmoidal trajectory, the growth
rate is at its maximum.

When no sign change is found within the evaluated interval, PGM reports
that no inflection point was detected.

Acceleration extrema

PGM defines:

(P\_1): time of maximum positive growth acceleration before the
inflection point; (P\_2): time of minimum, or most negative, growth
acceleration after the inflection point.

Thus,

\arg\max\_\{t\textless T\_i\} W'\,'(t), \$\$

and

\arg\min\_\{t\textgreater T\_i\} W'\,'(t). \$\$

These points describe the transition from increasing acceleration toward
maximum growth rate and the subsequent period of strongest deceleration.

Ontogenetic growth force

The ontogenetic growth force is defined as

\[
W'(t)W''(t).
\]

PGM may report:

(F\_1): maximum OGF before the inflection point; (F\_i): the
inflection-point location represented in the OGF sequence; (F\_2):
minimum OGF after the inflection point.

Because (W'(t)) is nonnegative for an increasing growth trajectory, the
sign of OGF generally follows the sign of (W'\,'(t)).

OGF-based points should be interpreted as geometric or operational
descriptors unless independent biological validation is available.

Effective growth onset

Available onset descriptors may include:

tangent-based onset; threshold-based onset; OGF-derived onset;
PAA-derived onset. Tangent method

The tangent method uses the tangent line at the point of maximum growth
rate and calculates its intersection with the baseline response
represented by the fitted value at the beginning of the evaluated
interval.

Its result depends on the observed or selected lower time boundary.

Threshold method

The threshold method identifies the first time at which the fitted
trajectory reaches a selected proportion of its fitted late-stage value.

The result depends directly on the chosen threshold.

OGF-derived onset

The OGF-derived onset is calculated using the third derivative of the
OGF before (P\_1).

It should be treated as a model-derived geometric descriptor rather than
as a universally established biological boundary.

PAA

PAA is identified before (P\_1) from a zero crossing of the fourth
derivative of the growth function.

It represents a transition in the evolution of growth acceleration.

Late-growth descriptors

Available late-stage descriptors may include:

an OGF-derived late point; PDA.

These are identified after (P\_2) using derivative-based rules.

They indicate geometric transitions toward the late portion of the
fitted trajectory and are not direct evidence that biological growth has
completely ceased.

Integral descriptors

PGM may calculate the area under the fitted growth curve:

\[
\int_{t_L}^{t_U} W(t),dt.
\]

It may also calculate the integrated squared growth rate:

\[
\int_{t_L}^{t_U}
\left[W'(t)\right]^2dt,
\]

and the response-weighted integrated squared growth rate:

\[
\int_{t_L}^{t_U}
W(t)\left[W'(t)\right]^2dt.
\]

These descriptors depend on the selected integration interval.

Comparisons require the same time unit, response unit, interval, and
fitted-variable definition.

Numerical evaluation

Critical points are located on a dense time grid spanning the fitted
interval.

Consequently, their reported precision is affected by the numerical grid
resolution and should not be interpreted as exact analytic precision.

Uncertainty

The current critical-point and integral tables provide point estimates.

Confidence intervals for derived times and integral descriptors are not
automatically implied by confidence intervals for the model parameters.

Bootstrap, simulation, or delta-method procedures would be required to
quantify their uncertainty separately.

\part{Academic use}

\chapter{reporting}\label{reporting}

Academic reporting Minimum information

An academic report should identify:

the PGM version; the jamovi version; the selected growth equation; the
equation parametrization; the response and time units; the number of
observations; the number of subjects, when longitudinal; the estimation
framework; the random-effects structure; the residual variance
structure; the residual correlation structure; the estimation method;
the model-comparison criteria; the residual diagnostics examined.
Ordinary nonlinear model

A possible methods statement is:

Growth trajectories were fitted using the PGM module for jamovi. The
selected nonlinear growth equation was estimated by nonlinear least
squares, with initial optimization performed using the
Levenberg--Marquardt algorithm. Model adequacy was evaluated using
parameter estimates, information criteria, prediction-error measures,
and residual diagnostics.

Generalized nonlinear least squares

When a variance function is included:

Growth trajectories were initially fitted by nonlinear least squares
using the Levenberg--Marquardt algorithm. The resulting estimates were
used to initialize a generalized nonlinear least-squares model with a
{[}power/exponential/constant-plus-power{]} residual variance function
based on the fitted response.

Nonlinear mixed-effects model

For longitudinal data:

Repeated growth measurements were analyzed using a nonlinear
mixed-effects model. Population growth parameters were represented by
fixed effects, and between-subject variation in the asymptotic parameter
was represented by a random effect. Parameters were estimated by maximum
likelihood after initialization using a Levenberg--Marquardt nonlinear
fit.

Residual autocorrelation

When applicable:

Within-subject residual dependence was modeled using a first-order
{[}AR(1)/CAR(1){]} correlation structure.

For CAR(1), the report may clarify:

Correlation was modeled as a function of the actual temporal distance
between observations.

Heteroscedasticity

For a power variance model:

Residual variance was modeled as a power function of the absolute fitted
response.

Avoid writing only that ``weighted regression'' was performed because
that does not identify the fitted variance model.

Model comparison

A possible statement is:

Candidate specifications were compared using AIC, AICc, and BIC,
together with residual diagnostics and biological plausibility of the
estimated parameters. Information criteria were compared only across
models fitted to the same response observations.

Derived growth descriptors

When reporting derivative-based descriptors:

Derived growth times were calculated from the fitted growth function and
its derivatives. These quantities were treated as model-based geometric
descriptors.

When uncertainty was not calculated, state:

Derived growth times are presented as point estimates; confidence
intervals were not calculated.

Tables

Parameter tables should report:

estimate; standard error; confidence interval; parameter unit; model
parametrization.

Model-comparison tables should report:

model; variance structure; correlation structure; random-effects
structure; number of observations; AIC; AICc; BIC; relevant error
metrics. Figures

Growth-curve figures should identify:

observed data; fitted trajectory; response and time units; group or
subject definition; critical-point labels, when shown.

Residual figures should be reported or summarized whenever model
assumptions are important to the conclusions.

Citation

The final software citation should include:

module authors; year; module title; version; software type; repository
or archival DOI.

The jamovi software and relevant statistical libraries should also be
cited according to the requirements of the target journal.

\chapter{Limitations}\label{limitations}

Limitations Model dependence

All parameter estimates and derived descriptors depend on the selected
mean function.

A good numerical fit does not guarantee that the equation is
biologically appropriate.

Different growth equations may produce similar fitted values but
different parameter estimates, asymptotes, inflection times, and
derivative-based descriptors.

Data coverage

Asymptotic parameters are weakly supported when observations do not
cover the late growth phase.

Early or truncated datasets may produce:

unstable asymptotes; strong parameter correlations; wide confidence
intervals; implausible extrapolation; convergence failures. Current
random-effects structure

The current nonlinear mixed-effects implementation includes a random
effect on the asymptotic parameter (A).

It does not automatically test all possible random-effects structures.

If individuals differ substantially in growth rate, inflection time, or
shape, the current structure may be insufficient.

Distributional assumptions

The fitted models rely on assumptions about:

residual normality; the selected variance function; random-effect
normality; the selected residual correlation structure.

Residual diagnostics can reveal important violations but cannot prove
that all assumptions are correct.

Numerical estimation

Nonlinear models may be sensitive to:

starting values; parameter scaling; sparse data; weakly identified
parameters; highly flexible equations; extreme observations.

Failure to converge should not be resolved merely by repeatedly changing
optimizer settings without reconsidering the data and model
specification.

Approximate confidence intervals

Parameter confidence intervals use asymptotic approximations and
first-order uncertainty propagation.

They may be inaccurate near parameter boundaries or under strong
nonlinearity.

Descriptive (R\^{}2)

The (R\^{}2) and adjusted (R\^{}2) values produced by PGM are
descriptive residual-based measures.

For mixed-effects models, they are not equivalent to specialized
marginal and conditional mixed-model (R\^{}2) measures.

Information criteria

Information criteria support relative comparison among compatible
models.

They do not:

test biological truth; establish residual adequacy; guarantee predictive
performance outside the observed range; replace subject-matter
interpretation. Critical points

Derivative-based critical points depend on:

the selected growth equation; the fitted parameters; the observed time
interval; numerical grid resolution; existence of the required extrema
or zero crossings.

A missing point may indicate that the mathematical condition was not
found within the evaluated interval.

It should not automatically be interpreted as evidence that the
biological process lacks the corresponding phase.

Uncertainty of derived descriptors

PGM currently reports point estimates for critical times and integral
descriptors.

Their uncertainty is not automatically quantified.

Extrapolation

Predictions beyond the observed time range should be treated cautiously,
especially when the asymptotic phase is not well represented in the
data.

\bookmarksetup{startatroot}

\chapter*{References}\label{references}
\addcontentsline{toc}{chapter}{References}

\markboth{References}{References}

\protect\phantomsection\label{refs}




\end{document}
